\documentclass[11pt]{article}
\usepackage[margin=1in]{geometry}
\usepackage{amsmath,amssymb,amsthm,hyperref,enumitem}
\newtheorem{theorem}{Theorem}
\newtheorem{lemma}{Lemma}
\newcommand{\BC}{\mathrm{BC}}
\newcommand{\Z}{\mathbb Z}
\title{A fourth-root pole in the seventh stained-glass column}
\author{Principia Math In Progress}
\date{18 September 2026}
\begin{document}
\maketitle

\paragraph{Task 1: resolve MathDB 347982.}
Blomberg, Shannon, and Sloane's Conjecture 4.3 predicts that each fixed
column \(k\ge3\) of their triangle, quadrilateral, and total-cell arrays
has a generating function with denominator
\[
 D_k(z)=(1-z^{k-2})(1-z^{k-1})(1-z^k).
\]
The statement remains in the published article and subsequent book
chapter \cite{BSS}. We disprove it using the total-cell column \(k=7\).
The argument also proves rationality of every fixed total-cell column.

Take the rectangle \([0,n]\times[0,1]\), with every pair of boundary
points \((j,0),(j,1)\), \(0\le j\le n\), joined by a segment. Write
\(c_{n,k}\) for its number of cells in
\((k-1,k)\times(0,1)\), and put
\[
 C_k(z)=\sum_{n\ge k}c_{n,k}z^n .
\]
The integer vertical segments already belong to this arrangement, so
its cells are contained in the unit strips. This is exactly the array
in Table 4 of the article, or Table 5.4 of the book \cite{BSS}, and
OEIS A333288 \cite{OEIS}.

\begin{theorem}
The rational continuation of \(C_7\) satisfies
\[
 \lim_{z\to i}(1-z^4)C_7(z)=2i .
\]
In particular, \(C_7(z)\) cannot be a polynomial divided by \(D_7(z)\),
and Conjecture 4.3 is false.
\end{theorem}

\paragraph{Counting cells through consecutive chords.}
For \(n\ge k\), exactly
\[
 L_{n,k}=2k(n-k+1)
\]
segments pass through the open unit strip: their two horizontal endpoint
indices lie on opposite sides of the gap between \(k-1\) and \(k\).
If \(m_v\) segments meet at an interior intersection \(v\), then
\begin{equation}\label{cells}
 c_{n,k}=1+L_{n,k}+\sum_v(m_v-1).
\end{equation}
Indeed, insert the segments one at a time. A new segment with \(h\)
distinct old interior intersections creates \(h+1\) new cells.
At a point eventually incident with \(m_v\) segments, exactly the last
\(m_v-1\) insertions count that point.

Represent a chord by the pair \((A,B)\) of its bottom and top endpoint
indices. At an intersection \((u,v)\), these pairs are the integer
points of
\[
 (1-v)A+vB=u,\qquad 0\le A,B\le n .
\]
Their consecutive differences have a unique primitive direction
\((p,-r)\), with \(p,r\ge1\) and \(\gcd(p,r)=1\).
There are exactly \(m_v-1\) consecutive pairs. This familiar
point-line description also appears in Karlsson's treatment of
simple intersections \cite{Karlsson}; here we count all consecutive
pairs and retain the unit-strip restriction.

Every such pair can be written uniquely as
\[
 (A,B)=(a,r+c),\quad(a+p,c),
 \qquad a,c\ge0,\quad a+p\le n,\quad r+c\le n.
\]
Conversely, a primitive difference makes these two chords consecutive:
there is no intervening lattice point, while all intermediate points
of the line inside the square would represent available chords.
Their intersection is
\[
 \left(\frac{ar+pc+pr}{p+r},\,\frac{p}{p+r}\right).
\]
Consequently, if \(A_{n,k}(p,r)\) counts the nonnegative integer pairs
\((a,c)\) with \(a+p,r+c\le n\) and
\begin{equation}\label{strip}
 (k-1)(p+r)<ar+pc+pr<k(p+r),
\end{equation}
then
\begin{equation}\label{primitive}
 c_{n,k}=1+2k(n-k+1)+
 \sum_{\substack{p,r\ge1\\(p,r)=1}} A_{n,k}(p,r).
\end{equation}
Both inequalities are strict: intersections on the strip boundary
must not be counted as interior intersections.

\paragraph{Separating the finite part from the tail.}
Reflection across the horizontal midline interchanges \(p,r\) and
\(a,c\). Thus we may restrict to \(p\le r\), doubling when \(p<r\).
The sole primitive diagonal case is \(p=r=1\).
Since \eqref{strip} implies \(pr<k(p+r)\), we have \(p<2k\).
For \(p>k\), it also gives
\[
 r<\frac{kp}{p-k}.
\]
These directions are finite in number. For each fixed direction,
\eqref{strip} bounds both \(a,c\). Such a pair contributes
\[
 \frac{z^{\max(k,a+p,r+c)}}{1-z}
\]
to \(C_k\), so all these finite contributions are analytic at \(i\).
The same holds for any finite set of the remaining directions.
The initial term \(1+L_{n,k}\) has generating function
\begin{equation}\label{base}
 \frac{z^k(1+2k-z)}{(1-z)^2}.
\end{equation}

Put \(T=k+1\). It remains to consider, for \(1\le p\le k\),
\[
 r=pt+s,\qquad t\ge T,\quad
 s\in U_p:=\{1\le s\le p:\gcd(s,p)=1\}.
\]
This representation includes \(U_1=\{1\}\), and these directions all
have \(r>p\). In particular \(r/p>k\).
Writing \(j=p+a\), the strip condition is equivalent to
\begin{equation}\label{cbounds}
 k-1+\frac{(k-1-j)r}{p}<c<
 k+\frac{(k-j)r}{p}.
\end{equation}
If \(j\ge k+1\), the upper bound is negative. If \(j=k\), exactly
\(c=0,\ldots,k-1\) occur. If \(j\le k-1\), the integer bounds are
\[
 \ell=k+\left\lfloor\frac{(k-1-p-a)r}{p}\right\rfloor,
 \qquad
 h=k-1+\left\lceil\frac{(k-p-a)r}{p}\right\rceil .
\]
Here \(\ell\ge0\) and \(h\ge\ell\), since their defining open interval
has length \(1+r/p>k+1\). Also \(r+c>k\ge p+a\), so the activation
threshold is \(n=r+c\).
Summing first over \(n\), then over \(c\), gives
\[
 \sum_{c=\ell}^h\frac{z^{r+c}}{1-z}
 =\frac{z^{r+\ell}-z^{r+h+1}}{(1-z)^2}.
\]
The two exponents simplify to
\[
 r+\ell=(k-1-a)t+k+
 \left\lfloor\frac{(k-1-a)s}{p}\right\rfloor,\qquad
 r+h+1=(k-a)t+k+
 \left\lceil\frac{(k-a)s}{p}\right\rceil .
\]
For \(j=k\), the contribution is
\((1-z^k)z^{pt+s}/(1-z)^2\).

Summing the geometric series over \(t\ge T\) proves that
\begin{equation}\label{tail}
\begin{split}
 R_k(z)=\frac{2}{(1-z)^2}
 \sum_{p=1}^k\sum_{s\in U_p}\Bigg[
 &\sum_{a=0}^{k-p-1}
 \left\{
 \frac{z^{(k-1-a)T+k+\lfloor(k-1-a)s/p\rfloor}}{1-z^{k-1-a}}
 \right.\\[-2pt]
 &\hspace{35mm}\left.
 -\frac{z^{(k-a)T+k+\lceil(k-a)s/p\rceil}}{1-z^{k-a}}
 \right\}
 +\frac{(1-z^k)z^{pT+s}}{1-z^p}
 \Bigg]
\end{split}
\end{equation}
differs from \(C_k(z)\) by \eqref{base} plus a polynomial divided
by \(1-z\). Empty sums have value zero. The calculation is valid
coefficientwise as a formal power series, since each coefficient
counts a finite arrangement. It exhibits \(C_k\) as a rational function.
In particular \(C_7-R_7\) is analytic at \(i\).

\paragraph{The fourth-root contribution.}
Set \(k=7\), \(T=8\). All periods in \eqref{tail} lie between 1 and 7,
so at \(z=i\) only denominator \(1-z^4\) can vanish.
The positive terms with period 4 have \(a=2\), \(p\le4\); the negative
ones have \(a=3\), \(p\le3\). The final term has period 4 precisely
when \(p=4\). Since \(i^{4T}=1\), their combined numerator is
\begin{equation}\label{B}
\begin{split}
 B={}&\sum_{p=1}^4\sum_{s\in U_p}
 i^{7+\lfloor4s/p\rfloor}
 -\sum_{p=1}^3\sum_{s\in U_p}
 i^{7+\lceil4s/p\rceil}
 +(1-i^7)(i+i^3).
\end{split}
\end{equation}
For \(p=1,2\), the floor and ceiling coincide, so those terms cancel.
For \(p=3\), the contribution is
\[
 (i^8-i^9)+(i^9-i^{10})=2.
\]
For \(p=4\), the positive sum is \(i^8+i^{10}=0\);
the final term is zero as well. Hence \(B=2\), and
\[
 \lim_{z\to i}(1-z^4)C_7(z)
 =\frac{2B}{(1-i)^2}=2i\ne0.
\]
But
\[
 D_7(i)=(1-i^5)(1-i^6)(1-i^7)=4\ne0.
\]
A polynomial divided by \(D_7\) is analytic at \(i\), which proves the
theorem. Shifting the column's starting index or changing finitely
many initial coefficients does not remove this pole.

\paragraph{Reproducibility.}
The accompanying exact verifier checks the primitive formula against
direct rational intersection enumeration, the published array data,
and the finite-part-plus-tail decomposition. It constructs the full
rational function for column 7 without guessing a recurrence.
Its numerator has degree 74 and is coprime to
\[
 (1+z^2)(1-z^5)(1-z^6)(1-z^7).
\]
Polynomial division further gives
\[
 D_7(z)C_7(z)=P(z)+\frac{4z}{1+z^2},
 \qquad \deg P=72.
\]
The integer coefficient certificate is included in the package.
These computations support the argument; the short pole calculation
above already proves the counterexample without a numerical cutoff.
The proof and source audits were performed by the authoring agent;
no independent peer review or proof-assistant certification is claimed.

\paragraph{Task 2: determine the minimal denominators.}
The decomposition gives the universal upper denominator
\((1-z)^2\operatorname{lcm}_{1\le d\le k}(1-z^d)\)
for the total-cell column. Which cyclotomic factors cancel as a function
of \(k\), and what is the exact eventual quasipolynomial period?

\paragraph{Task 3: separate triangles and quadrilaterals.}
Can the consecutive-chord count be refined to determine the two cell
types separately? Determine their rational generating functions and
the individual contributions to the fourth-root pole in column 7.

\paragraph{Task 4: simplify the column formulas.}
Can the finite and primitive-residue sums be reorganized using
arithmetic functions to give short coefficient formulas for general
\(c_{n,k}\), including a useful bound on the start of quasipolynomial
behavior?

\paragraph{Task 5: compare fixed and moving columns.}
Determine the leading quadratic coefficient for fixed \(k\), and
study how the asymptotics change when \(k\) grows with \(n\), especially
near the central column. These are followup questions, not additional
claimed solutions.

\begin{thebibliography}{9}
\bibitem{BSS}
L. Blomberg, S. R. Shannon, and N. J. A. Sloane,
\emph{Graphical Enumeration and Stained Glass Windows, 1:
Rectangular Grids},
Integers \textbf{21A} (2021), A5, Conjecture 4.3, p.~13.
\href{https://math.colgate.edu/~integers/graham5/graham5.pdf}{Journal text}.
Also in \emph{Number Theory and Combinatorics} (2022), pp.~65--98,
Conjecture 4.3, p.~78,
\href{https://doi.org/10.1515/9783110754216-005}{doi:10.1515/9783110754216-005}.
Preprint \href{https://arxiv.org/abs/2009.07918}{arXiv:2009.07918v2}.
\bibitem{Karlsson}
J. Karlsson, \emph{Church Windows}, 23 June 2026, 9 pp.,
\href{https://oeis.org/A334701/a334701_2.pdf}{OEIS-hosted manuscript}.
\bibitem{OEIS}
The OEIS Foundation Inc.,
\href{https://oeis.org/A333288}{A333288}, including L. Blomberg's
80-row data table. Accessed 18 September 2026.
\end{thebibliography}
\end{document}
